\section{Background}
\label{sec:background}

\subsection{The Transformer Circuits Framework}
\citet{elhage2021mathematical} developed a mathematical framework for decomposing
transformer computations into named circuits operating on the residual stream.
Each attention head writes a low-rank update; the final prediction is determined by
the accumulated sum projected through the unembedding matrix.
This linearity enables direct logit attribution: each head's contribution to the
final prediction can be computed without re-running the model.

\subsection{Induction Heads}
\citet{olsson2022context} identified induction heads as the primary mechanism
underlying in-context learning.
An induction head at layer~$l$ attends to the token that follows the previous
occurrence of the current token, implementing a copy-and-complete pattern:
if $A \to B$ was seen earlier, the model predicts $B$ after the next $A$.

In two-layer attention-only models, the circuit consists of: (1) a
\emph{previous-token head} in layer~0 that copies each token's identity to the
next residual-stream position; and (2) an \emph{induction head} in layer~1 that
reads this information to look back at previous occurrences.

\subsection{Activation Patching}
Activation patching~\citep{conmy2023automated} establishes causal claims by replacing
one component's activation from a corrupted run with the corresponding activation from
a clean run, measuring recovery of the target behaviour.
We use second-half shuffling as the corruption, following \citet{wang2022interpretability}.

\subsection{Circuit Stability Under Fine-Tuning}
Prior work has studied catastrophic forgetting at the capability level but circuit-level
stability during fine-tuning has not been studied.
\citet{marks2024sparse} demonstrated that circuits can be identified and edited, but not
their dynamics under gradient descent.
This paper fills that gap for the induction circuit.
